% 原书 p.88–p.92（PDF p.96–p.100）
% 第 4 章 4.1 节末尾（例 4.1.1）与 4.2 节：稳定性定理

% >>> p096
\begin{wexam}{4.1.1}
  回忆例 3.2.3，我们由 Liouville 域的边界构造了切触流形。辛椭球 $E(a,b)$ 是一个标准例子。
  在 $(\R^{2n},\omega_{\mathrm{std}}=\dd\lambda_{\mathrm{std}})$ 中，Liouville 域
  $(W,\omega_{\mathrm{std}})$ 也称为\textbf{星形域}（star-shaped domain），按定义即其边界
  与 (3.8) 中定义的径向向量场横截。其边界 $\partial W$ 上的切触结构由
  $\xi=\ker\lambda_{\mathrm{std}}|_{\partial W}$ 给出。

  现在，给定一族光滑星形域 $\{(W_t,\omega_{\mathrm{std}})\}_{t\in[0,1]}$，其边界
  $\partial W_t$ 诱导出一族切触结构 $\xi_t:=\ker\lambda_{\mathrm{std}}|_{\partial W_t}$。
  定理 4.1.2 表明，

% >>> p097
  对任意 $t\in[0,1]$，$(\partial W_0,\xi_0)$ 与 $(\partial W_t,\xi_t)$ 切触同构。注意
  $\partial W_t$ 在拓扑上是奇维球面 $S^{2n-1}$，因此一族光滑星形域在 $S^{2n-1}$ 上
  诱导出一族彼此切触同构的切触结构。
  % [Tutor Note: 原文疑误——原文 \"contactomprhic\" 系 \"contactomorphic\" 的拼写错误，
  % 译文按正确含义处理。详见 qsg/notes/review_ch04.md。]

  更具体地说，辛椭球 $E(a_1,\cdots,a_n)$ 边界上的切触结构本质上是相同的。然而，不同辛
  椭球的 Reeb 动力学（例如由闭 Reeb 轨道给出的数据）差别很大。原因在于 Reeb 动力学也
  由切触 $1$-形式决定。

  上述讨论与注 3.2.3 密切相关：该注断言，Liouville 流形
  $(\R^{2n},\omega_{\mathrm{std}})$ 存在一个理想边界 $M$，拓扑上是 $S^{2n-1}$，且在
  $M$ 上有一个「与切触 $1$-形式无关」的典范切触结构 $\xi$。上面那些彼此切触同构的
  切触结构 $\xi_t$，正是通过投影 $\R^{2n}\to M$ 的不同截面来实现的。\hfill$//$
\end{wexam}
% [Tutor Note: 术语依据——原文 \"contact 1-form free\" 应读作「不依赖切触 1-形式的选取」
% （即由 ξ 本身决定的典范结构），而不是「没有切触 1-形式的」。按后者理解会与切触结构的
% 定义矛盾：任何余维 1 切触结构都局部地是某个切触 1-形式的核。]

\section{稳定性定理（Stability Theorems）}

\begin{keytakeaways}
\begin{itemize}
  \item 几何图像：不止「恰当扰动」可用 Moser 论证；只要辛形式族的 de Rham 上同调类不变，
        就能用同痕把它们拉回到同一个辛形式。二维闭曲面情形完全由亏格与面积决定。
  \item 数学工具：Hodge 理论（Hodge $\ast$ 算子、余微分 $\delta$、Laplace 算子与 Hodge
        分解），以及 $L^2$ 配对
        $\langle\alpha,\beta\rangle=\int_M\alpha\wedge\ast\beta$。
  \item 本节结论：Moser 体积形式定理（定理 4.2.1），以及「上同调类不变即可同痕」的
        定理 4.2.2。
\end{itemize}
\end{keytakeaways}

要应用定理 4.1.1，就需要验证条件 (4.1)。人们常遇到的情形来自 $\omega_0$ 的一个
「线性」形变，即
\begin{equation}\tag{4.7}
  \omega_t=\omega_0+t\,\dd\sigma
  \quad\text{其中}\ \sigma\in\Omega^1(M).
\end{equation}
我们要强调定理 4.1.1 中一个隐含的条件：(4.7) 中的 $\sigma$ 要取得使 $\omega_t$ 对一切
$t\in[0,1]$ 都是辛的（特别地，非退化）（参见例 2.3.1）。
% [Tutor Note: 易错预警——「线性形变」未必自动保持辛性：ω_0 与 ω_1 都辛，并不能保证
% 中间的 ω_t = ω_0 + t dσ 都对一切 t 非退化（参见注 4.2.1 所引 [150] 中定理 2.1）。
% 此外，(4.7) 中的 σ 与 (4.1) 中的 σ_t 是同一角色：取 σ_t ≡ σ 即满足 dω_t/dt = dσ。]

这种线性形变 (4.7) 在定向闭曲面（即二维辛流形）上非常有效。下面用例子来说明。

\begin{wexam}{4.2.1}
  设 $M$ 是闭定向光滑曲面，$\omega_0,\omega_1$ 是 $M$ 上的辛形式。则存在 $M$ 上的
  微分同胚 $\varphi$ 使得 $\varphi^*\omega_1=\omega_0$，当且仅当二者具有相同的积分：
  \begin{equation}\tag{4.8}
    \int_M\omega_0=\int_M\omega_1.
  \end{equation}
  事实上，若这样的 $\varphi$ 存在，则对 $\varphi^*\omega_1=\omega_0$ 两边积分，得
  \begin{equation*}
    \int_M\omega_0=\int_M\varphi^*\omega_1=\int_{\varphi(M)}\omega_1=\int_M\omega_1.
  \end{equation*}
  反之，设 (4.8) 成立，则 $[\omega_1-\omega_0]=0\in H^2(M;\R)$。于是存在
  $\beta\in H^1(M;\R)$，使得 $\omega_1=\omega_0+\dd\beta$。考虑线性组合
  \begin{equation*}
    \omega_t=\omega_0+t\,\dd\beta\quad\text{对}\ t\in[0,1].
  \end{equation*}
  所得的 $\{\omega_t\}_{t\in[0,1]}$ 是一族光滑非退化体积形式。于是定理 4.1.1 适用，
  我们得到 $M$ 上的同痕 $\{\varphi_t\}$，使得对任意 $t\in[0,1]$ 都有
  $\varphi_t^*\omega_t=\omega_0$。特别地，$\varphi=\varphi_1$ 满足
  $\varphi^*\omega_1=\omega_0$。\hfill$//$
\end{wexam}
% [Tutor Note: 原文疑误——此处 β ∈ H^1(M;R) 应为 β ∈ Ω^1(M)：H^1(M;R) 的元素是上同调类，
% 不能直接写 dβ；正确的表述是「存在 1-形式 β 使得 ω_1 = ω_0 + dβ」。译文按原文照译，
% 详见 qsg/notes/review_ch04.md。]

% >>> p098
下面这个练习从辛几何的视角给出闭定向曲面的分类。事实上，在二维情形，辛结构就是体积形式，
而这些曲面之间的辛同胚就是保体积且保定向的微分同胚。在如此强的要求下，下面的练习无非是说：
分类完全由亏格与面积（体积）决定。
% [Tutor Note: 含糊点拨——练习 4.2.1 的证明要点：由曲面分类定理先取微分同胚，再用
% 例 4.2.1 的 Moser 论证把两个面积形式同痕化。二维时闭 2-形式自动非退化（面积形式），
% 且同上同调类的两形式之差恰当，故线性形变全程非退化；这是例 4.2.1 能直接适用的原因。]

\begin{wexer}{4.2.1}
  证明：二维闭辛流形 $(M_0,\omega_0)$ 与 $(M_1,\omega_1)$ 辛同构，当且仅当它们具有相同的
  亏格与面积。这里的面积由 $\int_{M_i}\omega_i$ 计算，$i=0,1$。
\end{wexer}

同样地，例 4.2.1 中的论证对一般闭流形与体积形式也成立。确切地说，我们有下面结果。

\begin{wthm}{4.2.1}（Moser，[164] 中的定理 2）
  设 $M$ 是闭定向流形，其上装有两个体积形式 $\tau_0$ 与 $\tau_1$。则存在 $M$ 上的
  微分同胚 $\varphi$ 使得 $\varphi^*\tau_1=\tau_0$（即保体积微分同胚），当且仅当
  $\tau_0$ 与 $\tau_1$ 在 $M$ 上具有相同的体积积分。
\end{wthm}

下面的结果是 Moser 论证的又一应用，其充分条件比条件 (4.1) 更一般。

\begin{wthm}{4.2.2}
  设 $M$ 是闭流形，$\{\omega_t\}$ 是 $M$ 上的一族光滑辛结构，且 $[\omega_t]$ 在
  $H^2(M;\R)$ 中代表同一个上同调类。则存在 $M$ 上的光滑同痕
  $\{\varphi_t\}_{t\in[0,1]}$，使得 $\varphi_t^*\omega_t=\omega_0$。
\end{wthm}

在给出证明之前，需要一些 Hodge 理论的背景（详细的介绍见 [239] 第 4 章）。回忆：在
$2n$ 维流形 $M$ 上，\textbf{Hodge $\ast$ 算子}是一个映射
$\ast\colon\Omega^k(M)\to\Omega^{2n-k}(M)$，它把任意 $k$-形式
$\beta\in\Omega^k(M)$ 通过如下非退化配对映为唯一的 $(2n-k)$-形式（记作
$\ast\beta\in\Omega^{2n-k}$）：
\begin{equation}\tag{4.9}
  \langle\alpha,\beta\rangle:=\int_M\alpha\wedge\ast\beta,
\end{equation}
其中 $\alpha$ 是任意 $k$-形式。这里 $\langle\alpha,\beta\rangle$ 给出由 Riemann 度量
诱导的 $\alpha$ 与 $\beta$ 的内积。换言之，这个配对给出形式空间上的 $L^2$-范数。在此
配对下，外微分 $\dd_*\colon\Omega^*(M)\to\Omega^{*+1}(M)$ 具有伴随算子
$\delta_*\colon\Omega^*(M)\to\Omega^{*-1}(M)$，满足如下关系
\begin{equation*}
  \langle\dd\omega,\eta\rangle=\langle\omega,\delta\eta\rangle,
\end{equation*}
其中 $\omega\in\Omega^*(M)$，$\eta\in\Omega^{k+1}(M)$。定义 Laplace 算子为
$\Delta_*:=\dd_*\delta_*+\delta_*\dd_*\colon\Omega^*(M)\to\Omega^*(M)$。对任意
$k\in\N$，记
\begin{equation}\tag{4.10}
  \mathcal H_\Delta^k(M):=\ker\bigl(\Delta|_{\Omega^k(M)}\bigr)
  =\{\alpha\in\Omega^k(M)\mid \dd_k\alpha=\delta_k\alpha=0\}
\end{equation}
% >>> p099
为 $M$ 上的调和形式之集。于是 Hodge 分解定理断言：关于 $L^2$-范数 (4.9)，有如下正交
分解
\begin{equation}\tag{4.11}
  \Omega^k(M)=\operatorname{im}(\dd_{k-1})\oplus\operatorname{im}(\delta_{k+1})
  \oplus\mathcal H_\Delta^k(M)
\end{equation}
对任意 $k\in\N$ 成立。特别地，限制
$\Delta|_{\operatorname{im}(\dd)\oplus\operatorname{im}(\delta)}$ 是可逆的。此外，
逆映射
$(\Delta|_{\operatorname{im}(\dd)\oplus\operatorname{im}(\delta)})^{-1}$
在光滑性意义下保持正则性（见 [231] 第 1.5 节）。
% [Tutor Note: 含糊点拨——(4.11) 的「正交分解」与「Δ 在 im(d)⊕im(δ) 上可逆」依赖三条
% 经典事实：(a) 紧流形上 Hodge 分解成立（需 Riemann 度量与椭圆正则性）；(b) Δ 与 d、δ
% 交换；(c) Hodge 调和形式在每阶有限维。逆算子 Δ^{-1} 保持光滑性即椭圆正则性定理，
% 这是 (4.13) 中 σ_t 光滑的关键，不能省略。参考 [239] 第 4 章、[231] 第 1.5 节。]

\begin{wproof}[定理 4.2.2 的证明（朴素尝试）]
  对任意 $t\in[0,1]$，$\omega_t-\omega_0$ 是恰当的（因为由假设
  $[\omega_t-\omega_0]=[\omega_t]-[\omega_0]=0$），所以存在 $1$-形式 $\alpha_t$
  使得 $\omega_t=\omega_0+\dd\alpha_t$。然而 $\alpha_t$ 只是对每个 $t$ 逐点定义的，
  这不能保证 $1$-形式的道路 $\{\alpha_t\}_{t\in[0,1]}$ 可微（即使
  $\{\omega_t\}_{t\in[0,1]}$ 光滑）。或者，注意
  \begin{equation}\tag{4.12}
    \Bigl[\frac{\dd}{\dd t}\omega_t\Bigr]=\frac{\dd}{\dd t}[\omega_t]=0.
  \end{equation}
  事实上，(4.12) 中的第二个等号来自我们的假设，而第一个等号来自广义 Leibniz 法则：
  对任意 $2$-闭链 $C$，有
  \begin{equation*}
    \frac{\dd}{\dd t}\int_C\omega_t=\int_C\frac{\dd}{\dd t}\omega_t.
  \end{equation*}
  于是对每个 $t\in[0,1]$，存在 $1$-形式 $\beta_t$ 使得
  $\dd\omega_t/\dd t=\dd\beta_t$。同样，这不能保证 $1$-形式的道路
  $\{\beta_t\}_{t\in[0,1]}$ 可微。
  % [Tutor Note: 含糊点拨——朴素尝试的困难不在「存在性」，而在「光滑依赖于 t」：
  % 逐点取原像（或逐点解 dβ_t = dω_t/dt）可以用线性代数或 Poincaré 引理完成，但如此得到
  % 的 t ↦ α_t 未必光滑。作者随后用 Hodge 理论给出一族在 t 上光滑的典范选择，这才是
  % 「正确进路」的实质。]

  \textbf{（正确进路 I）}为得到 (4.1) 中所需的光滑族
  $\{\sigma_t\}_{t\in[0,1]}$，Hodge 理论可以帮忙。由于 $\dd\omega_t/\dd t$ 恰当，
  它当然落在 $\operatorname{im}(\dd)\oplus\operatorname{im}(\delta)$ 中，而 Laplace
  算子在其上可逆。考虑
  \begin{equation}\tag{4.13}
    \sigma_t:=\delta\Bigl(\Delta^{-1}\Bigl(\frac{\dd\omega_t}{\dd t}\Bigr)\Bigr)
    \in\Omega^1(M).
  \end{equation}
  那么 $\{\omega_t\}_{t\in[0,1]}$ 的光滑性蕴含 $\{\sigma_t\}_{t\in[0,1]}$ 的光滑性。
  此外，我们断言由 (4.13) 定义的 $\sigma_t$ 满足
  \begin{equation}\tag{4.14}
    \frac{\dd}{\dd t}\omega_t=\dd\sigma_t.
  \end{equation}
  于是定理 4.1.1 可以应用。为验证条件 (4.14)，记
  $\tau_t:=\Delta^{-1}(\dd\omega_t/\dd t)$。则 (4.14) 可改写为
  $\Delta\tau_t=\dd(\delta\tau_t)$，这等价于 $\delta(\dd\tau_t)=0$。为此，我们
  将求助于非退化配对 (4.9)。具体地说，考虑 $\delta(\dd\tau_t)$ 的 $L^2$-范数，有
  \begin{equation*}
    \begin{aligned}
      \bigl\langle\delta(\dd\tau_t),\delta(\dd\tau_t)\bigr\rangle
      &=\Bigl\langle\frac{\dd}{\dd t}\omega_t-\dd(\delta\tau_t),
      \delta(\dd\tau_t)\Bigr\rangle\\
      &=\Bigl\langle\dd\Bigl(\frac{\dd}{\dd t}\omega_t-\dd(\delta\tau_t)\Bigr),
      \dd\tau_t\Bigr\rangle
    \end{aligned}
  \end{equation*}
% >>> p100
  \begin{equation*}
    =\Bigl\langle\frac{\dd}{\dd t}(\dd\omega_t),\dd\tau_t\Bigr\rangle=0,
  \end{equation*}
  这说明 $\delta(\dd\tau_t)=0$。证毕。

  \textbf{（正确进路 II）}对每个 $t\in[0,1]$，由于 $\dd\omega_t/\dd t$ 恰当，存在
  $1$-形式 $\beta_t$ 使得
  $\dd\omega_t/\dd t=\dd\beta_t\in\operatorname{im}(\dd)\oplus
  \operatorname{im}(\delta)$。由 (4.11)，把逆 $\Delta^{-1}(\dd\beta_t)$ 写成
  $\dd a_t+\delta b_t+c_t\in\Omega^2(M)$，其中 $a_t,b_t$ 是某些形式，
  $c_t\in\mathcal H_\Delta^2(M)$。于是有
  \begin{equation*}
    \dd(\delta(\dd a_t))+\delta(\dd(\delta b_t))=\dd\beta_t,
  \end{equation*}
  特别地，$\delta(\dd(\delta b_t))\in\operatorname{im}(\dd)$。再次由 (4.11)，
  $\operatorname{im}(\dd)\cap\operatorname{im}(\delta)=\{0\}$，这蕴含
  $\delta(\dd(\delta b_t))=0$（从而 $\dd\beta_t=\dd(\delta(\dd a_t))$）。考虑
  $\sigma_t:=\delta(\Delta^{-1}(\dd\beta_t))$，则
  \begin{equation*}
    \dd\sigma_t=\dd(\delta(\dd a_t+\delta b_t+c_t))
    =\dd(\delta(\dd a_t))=\dd\beta_t=\frac{\dd}{\dd t}\omega_t,
  \end{equation*}
  其中第二个等号来自 $\delta\circ\delta=0$ 以及由 (4.10) 得到的 $\delta c_t=0$。
  证毕。
  % [Tutor Note: 易错预警——「im(d) ∩ im(δ) = {0}」是由 (4.11) 中直和分解在「同一阶」
  % 上的正交性得到的，应用时须核对次数：这里 δ(d(δb_t)) 是 2-形式，既在 im(δ) 中又在
  % im(d) 中，故为零；而 σ_t 是 1-形式，故 dσ_t 是 2-形式。若把 Δ^{-1} 与 d、δ 的
  % 交换次序搞错，就会得到错误的次数或符号。]
\end{wproof}

\begin{wrem}{4.2.1}
  在定理 4.2.2 中，如果对于给定的（光滑）辛结构道路 $\{\omega_t\}_{t\in[0,1]}$ 只知道
  $[\omega_0]=[\omega_1]$，则定理 4.2.2 的结论可能不成立，这已由 [150] 中的定理 2.1
  说明。
\end{wrem}
% [Tutor Note: 易错预警——注 4.2.1 说明「端点同调类相同」不足以保证可同痕，必须要求
% 整条道路的每个 ω_t 都与 ω_0 同调类一致（定理 4.2.2 的假设）。[150] 中的反例正是
% 一条中途离开了给定上同调类的辛道路；这是 Moser 论证的典型适用范围边界。]
